% ---
% id: meat-core/frank-wolfe/rem:fo-check-exposures
% kind: rem
% title: "One-sided exposures of the runtime checks"
% summary: >-
%   The rank comparison for (H-SUPP) and the three iterate checks are computed
%   in binary64 and leave three one-sided exposures, two on the unsafe side,
%   that no claim of this module bounds. A reference rank computed too large is
%   harmless, and the acceptance factor, exact by construction, adds none.
% prominence: standard
% status: accepted
% uses: [thm:firstorder, rem:fo-danger, rem:acc-side]
% mentions: [rem:fo-one-ruler, rem:fo-iterate-checks, req:verified-feasible-side, rem:fo-arithmetic]
% assumes: []
% refs: []
% provenance: newly-derived
% lean:
% ---

\begin{remark}\label{rem:fo-check-exposures}
The rank comparison of Remark~\ref{rem:fo-one-ruler} buys a refusal on the
safe side in the common case, a rank-deficient block being refused rather
than differentiated. It does not buy the hypothesis, and there are
\emph{three} one-sided residual exposures, the first two on the unsafe side
of Remark~\ref{rem:fo-danger} and the third on the membership hypothesis.
\emph{First}, a block whose true support is a proper subspace of
$W_{\hat c,s}$ but whose smallest relevant eigenvalue is reported above the
cutoff by the arithmetic of the spectral decomposition is differentiated on
a non-maximal face silently. \emph{Second}, a $\dim W_{\hat c,s}$ computed
\emph{too small} --- a direction of $W_{\hat c,s}$ that is genuinely present
but whose relative weight in
$\mathbf K_{\hat c,s}\mathbf K_{\hat c,s}^{\dagger}$, $\sigma^2$ over the
largest $\sigma^2$, is at most the cutoff, and which is therefore not
counted --- lowers the bar in \eqref{eq:hsupp-check} and passes a
rank-deficient block for the same reason. A $\dim W_{\hat c,s}$ computed too
\emph{large} is harmless for the claim: it can only fail the comparison,
which removes the iterate's bracket from $\hlb$ --- at every iterate, so
that $\hlb$ is withdrawn, if it counts a direction no
$\operatorname{rank}X_{\hat c,s}$ can count, which is why the two sides are
read with one ruler. \emph{Third}, and belonging to the checks of
Remark~\ref{rem:fo-iterate-checks} rather than to the rank comparison: each
of the three quantities those checks compare is itself computed in binary64.
$\lmin$ of an iterate block is computed, so a block whose true minimum
eigenvalue is negative within the eigensolver's own backward error passes at
exactly $0$ and is not repaired, and the reassembly of a block that
\emph{is} repaired is itself rounded; and the coupling residual
$\|\sum_s\tr_{\Qsys_s}J_s-\idon{A_0}\|_{\max}$ and each $\bowtie_k$ residual
are formed by partial traces and sums that round, so a point whose
\emph{true} residual exceeds its window can be measured inside it. So the
consumed point lies in $\Feas^{J}(\varepsilon)$ by construction only up to
that arithmetic.

\emph{The acceptance factor adds no fourth exposure to this list}: it is not
measured at all, being exact by construction on the \tref{exact island}
(Remark~\ref{rem:acc-side}), so the count is three because the three are the
$J$-side checks, and the product's other factor is established in a way
that has no residual to expose. Nor does the feasible side of a bracket:
its value is a certification computed exactly from the numbers it is given
(Requirement~\ref{req:verified-feasible-side}) and decides nothing at a
window, the rounding of the frame its rows are built with being charged
within it, and how far the gradient it is given lies from the exact one is
the binary64 evaluation of Remark~\ref{rem:fo-arithmetic}, not a check. No
claim of this module bounds any of the
three errors; what the declared thresholds and the exact-$0$ check buy is
that each comparison is well posed, that its reference is fixed by this
module rather than chosen by an implementation, and that no window
\emph{this module declares} admits a point at which the theorem does not
apply.
\end{remark}
